\chapter{Các Công Trình Liên Quan}

\section{Nghiên cứu về cây quyết định}
Cây quyết định (Decision Tree – DT) là một trong những mô hình học máy cơ bản, được ưa chuộng nhờ tính minh bạch và dễ hiểu. Các thuật toán kinh điển như \textbf{CART} \parencite{breiman1984cart} và \textbf{ID3} \parencite{quinlan1986induction} xây dựng cây theo hướng \emph{tham lam} (greedy): tại mỗi nút, đặc trưng được chọn dựa trên tiêu chí cục bộ (Gini impurity cho CART, information gain cho ID3), sau đó lá được gán nhãn theo đa số lớp hoặc quy tắc đơn giản khác.

Các nghiên cứu gần đây mở rộng cây quyết định cho các mục tiêu phi truyền thống, chẳng hạn \emph{optimal trees} \parencite{hu2019optimal, aglin2020learning} và tìm kiếm ngẫu nhiên (stochastic search). Những công trình này chứng minh DT có thể được điều chỉnh không chỉ để phân loại mà còn để giải các bài toán tối ưu có cấu trúc, đồng thời vẫn duy trì tính minh bạch.

\section{So sánh các hướng tiếp cận trước đó}

\subsection{Tiêu chí đánh giá}
Các phương pháp giải thích phản thực được so sánh theo bốn tiêu chí:

\begin{itemize}
\item \textbf{Mục tiêu tối ưu:} cục bộ (local) cho từng cá thể hay toàn cục (global) cho toàn bộ dữ liệu.
\item \textbf{Phủ toàn cục:} mức độ đảm bảo mọi cá thể đều có hành động/lý do cụ thể.
\item \textbf{Minh bạch:} khả năng diễn giải rõ ràng, dễ hiểu, và có lý do đi kèm.
\item \textbf{Nhất quán:} tính thống nhất trong việc gán hành động cho các cá thể tương tự.
\end{itemize}

\begin{table}[H]
\centering
\renewcommand{\arraystretch}{1.3}
\caption{So sánh các hướng tiếp cận}
\resizebox{\textwidth}{!}{%
\begin{tabular}{|>{\centering\arraybackslash}m{6cm}|>{\centering\arraybackslash}m{2cm}|>{\centering\arraybackslash}m{4cm}|>{\centering\arraybackslash}m{2.5cm}|>{\centering\arraybackslash}m{2.5cm}|}
\hline
Phương pháp & Mục tiêu tối ưu & Phủ toàn cục & Minh bạch & Nhất quán \\
\hline
CE \parencite{wachter2018counterfactualexplanationsopeningblack, karimi2021surveyalgorithmicrecoursedefinitions, Kanamori:AISTATS2022, schut2021} & Local & Không & Có & Có \\
CE \parencite{Ustun_2019, mothilal} & Local & Không & Có & Có \\
AReS \parencite{rawal2020ares} & Global & Không & Không & Không \\
MAME \parencite{ramamurthy2020} & Global & Có & Có & Không \\
GIME \parencite{gao2021learning} & Global & Có & Có & Không \\
\textbf{CET (Counterfactual Explanation Trees)} & Global & Có & Có & Có \\
\hline
\end{tabular}}
\end{table}

\subsection{Phân tích chi tiết}

\begin{enumerate}
\item \textbf{CE cục bộ} \parencite{wachter2018counterfactualexplanationsopeningblack, karimi2021surveyalgorithmicrecoursedefinitions, Kanamori:AISTATS2022, schut2021}:
Tối ưu chi phí để tìm hành động cho một cá thể duy nhất. Minh bạch và nhất quán ở cấp cá thể, nhưng không bao phủ toàn cục.

\item \textbf{CE cục bộ nâng cao} \parencite{Ustun_2019, mothilal}: 
Sử dụng \textbf{MILO/MILP} để sinh hành động khả thi với ràng buộc ngân sách và hướng thay đổi. (DiCE) sinh nhiều hành động đa dạng. Tuy nhiên, cả hai vẫn chỉ ở mức local, chưa bao phủ toàn cục.  

\item \textbf{AReS} \parencite{rawal2020ares}:  
Học tập hợp luật hai tầng để cung cấp recourse toàn cục. Do giới hạn số và độ dài luật, nhiều điểm không được bao phủ; các luật có thể chồng lấn, dẫn đến thiếu minh bạch và nhất quán.  

\item \textbf{MAME} \parencite{ramamurthy2020}:  
Tổng hợp giải thích cục bộ bằng phân cụm phân cấp để tạo mẫu giải thích chung. Đảm bảo phủ toàn cục và dễ hiểu, nhưng các cụm có thể chồng lấn, khiến một điểm dữ liệu có nhiều hành động khác nhau.  

\item \textbf{GIME} \parencite{gao2021learning}:  
Sử dụng \textit{topic modeling} để nhóm đặc trưng thành chủ đề và giải thích dựa trên các chủ đề đó. Minh bạch và bao phủ toàn cục, song một điểm có thể liên quan đến nhiều chủ đề, gây thiếu nhất quán.  

\item \textbf{CET (Counterfactual Explanation Trees):}  
Phân hoạch toàn bộ không gian đầu vào bằng cây quyết định. Mỗi cá thể thuộc đúng một lá, được gán duy nhất một hành động–lý do, đảm bảo minh bạch, nhất quán và phủ toàn cục.  


\end{enumerate}

\paragraph{Nhận xét:} So với CE cục bộ và các framework tổng hợp toàn cục (AReS, MAME, GIME), \textbf{CET} là phương pháp duy nhất đồng thời đảm bảo \textbf{minh bạch, nhất quán và phủ toàn cục}.

\section{Khoảng trống nghiên cứu}
CE cục bộ \parencite{wachter2018counterfactualexplanationsopeningblack, karimi2021surveyalgorithmicrecoursedefinitions, Ustun_2019} mang lại giải thích chi tiết cho từng cá thể nhưng thiếu khả năng bao phủ toàn cục. Trong khi đó, các framework toàn cục như AReS, MAME và GIME cung cấp góc nhìn tổng thể nhưng thường thiếu minh bạch hoặc nhất quán.

Do đó, khoảng trống nghiên cứu nằm ở việc phát triển một phương pháp vừa đảm bảo \textbf{toàn cục}, vừa duy trì \textbf{minh bạch và nhất quán}. CET được đề xuất để lấp đầy khoảng trống này thông qua cấu trúc cây quyết định kết hợp tối ưu hóa toàn cục.

\section{Vị trí nghiên cứu trong bối cảnh rộng hơn}
CET khắc phục đồng thời hạn chế của phương pháp cục bộ và toàn cục bằng hai ý tưởng chính:
\begin{enumerate}
\item \textbf{Phân hoạch không gian đầu vào:} mỗi đường đi từ gốc đến lá định nghĩa một luật rõ ràng, giúp lý giải trực quan.
\item \textbf{Tối ưu hành động theo vùng:} mọi điểm trong cùng một lá nhận cùng một hành động duy nhất, đảm bảo nhất quán và phủ toàn cục.
\end{enumerate}


\paragraph{Minh bạch:} Quy tắc if-then-else trong cây thể hiện rõ lý do cho từng hành động.
\paragraph{Nhất quán:} Mỗi điểm dữ liệu thuộc đúng một lá, do đó chỉ có một hành động duy nhất.
\paragraph{Phủ toàn cục:} Cây phân hoạch toàn bộ không gian, mọi cá thể đều có hành động gợi ý.

\paragraph{Kết luận:} CET là sự kết hợp giữa tính diễn giải của cây quyết định và tính linh hoạt của các phương pháp giải thích phản thực hiện đại, tạo thành một framework toàn diện cho giải thích vừa minh bạch vừa khả thi.


